\documentclass[11pt,a4paper]{article}

\usepackage[T1]{fontenc}
\usepackage{lmodern}
\usepackage[margin=2.3cm]{geometry}
\usepackage{amsmath,amssymb}
\usepackage{booktabs,tabularx,array}
\usepackage{enumitem}
\usepackage[dvipsnames]{xcolor}
\usepackage[colorlinks=true,linkcolor=MidnightBlue,urlcolor=MidnightBlue,citecolor=MidnightBlue]{hyperref}
\usepackage{microtype}

\setlist{itemsep=2pt,topsep=4pt}
\renewcommand{\arraystretch}{1.25}
\newcolumntype{L}{>{\raggedright\arraybackslash}X}

\newcommand{\code}[1]{\texttt{\small #1}}
\newcommand{\Naux}{N_{\mathrm{aux}}}
\newcommand{\nt}{n_{t}}
\newcommand{\Np}{N_{p}}
\newcommand{\Nfrag}{N_{\mathrm{frag}}}
\newcommand{\ov}{ov}
\newcommand{\Mt}{\tilde{M}}
\newcommand{\Tr}{\operatorname{Tr}}
\newcommand{\diag}{\operatorname{diag}}

\title{RIRPA in Vayesta\\[2pt]
\large Current state, cost, and a plan for HHT/CDERI improvements}
\author{}
\date{10 November 2026}

\begin{document}
\maketitle

\noindent
Nothing in this note was run. All scaling statements are operation counts. The measured
numbers in Section~\ref{sec:established} are quoted from the \code{monomial\_GW} and
\code{Davidson-gw} records, and the source is named each time.

\begin{table}[h]
\centering
\begin{tabularx}{0.9\textwidth}{@{}lL@{}}
\toprule
Symbol & Meaning \\
\midrule
$o$, $v$ & occupied / virtual orbitals of the full system \\
$\ov$ & number of spatial particle--hole pairs, $o\,v$ (spin-orbital codes use $2\ov$) \\
$\Naux$ & number of auxiliary (density-fitting) functions; $\mathcal{O}(N)$ \\
$L$ & occupied--virtual CDERIs, shape $(\Naux, \ov)$ \\
$D$ & diagonal of orbital-energy differences $\epsilon_a-\epsilon_i$, length $\ov$ \\
$\nt$ & number of rows of \code{target\_rot}, the space the moment is projected onto \\
$\Np$ & number of quadrature points (Clenshaw--Curtis) or poles (HHT) \\
$\Nfrag$ & number of fragments in an embedding calculation \\
\bottomrule
\end{tabularx}
\end{table}

%==========================================================================
\section{Vayesta as a whole}

Vayesta is a Python quantum-embedding package built on PySCF. It partitions a molecule or
solid into fragments, builds a cluster for each fragment (fragment orbitals plus a bath),
solves each cluster with a correlated solver (CCSD, FCI, MP2, EBFCI, \dots), and assembles
global properties from the fragment results. The main subpackages are:
\begin{itemize}
  \item \code{core/}: fragmentation, bath construction (\code{core/bath/}: DMET, MP2 and RPA
        bath natural orbitals, bosonic baths), integrals (\code{core/eris.py}, where
        \code{get\_cderi} lives), screening (\code{core/screening/}), and the generic embedding
        driver \code{core/qemb/}.
  \item \code{ewf/}, \code{dmet/}, \code{edmet/}: the embedding methods themselves
        (wave-function embedding, DMET, and extended/bosonic DMET).
  \item \code{solver/}: cluster solvers, including the electron--boson ones that EDMET needs.
  \item \code{rpa/}: full-system RPA. \code{rpa.py}, \code{ssrpa.py} and \code{ssurpa.py} are
        dense $\mathcal{O}(N^6)$ reference implementations. \code{rirpa/} is the
        density-fitted, moment-based code that the embedding layers use in production.
\end{itemize}

RIRPA matters because it supplies \emph{global long-range information to local clusters}.
It is called from four places:

\begin{center}
\small
\begin{tabularx}{\textwidth}{@{}>{\raggedright\arraybackslash}p{4.2cm}LL@{}}
\toprule
Caller & What it asks RIRPA for & \code{target\_rot} \\
\midrule
\code{core/screening/ screening\_moment.py} (\code{mrpa} screening, cRPA)
  & zeroth dd moment, then only $T_f\,\eta_0\,T_f^{\mathsf T}$ per fragment
  & stacked cluster ph spaces, spin-orbital, $\nt=\sum_f \ov_f^{\mathrm{clus}}$ \\
\code{core/bath/rpa.py} (RPA bath NOs)
  & zeroth moment, then $m_0 T^{\mathsf T}$
  & $R_o\otimes R_v$: $o\,v_{\mathrm{clus}}$ or $o_{\mathrm{clus}}\,v$ rows; \textbf{one RIRPA object per fragment} \\
\code{core/qemb/qemb.py} (bosonic bath; external dRPA correction)
  & zeroth moment / total energy
  & bosonic targets; or $L$ itself (cumulant energy) \\
\code{edmet/edmet.py}
  & zeroth moment and energy \textbf{with an xc kernel} (\code{rixc})
  & cluster ph and bosonic spaces \\
\bottomrule
\end{tabularx}
\end{center}

%==========================================================================
\section{How \texttt{vayesta/rpa/rirpa} is implemented}

\subsection{Files and classes}

\begin{center}
\small
\begin{tabularx}{\textwidth}{@{}>{\raggedright\arraybackslash}p{3.2cm}L@{}}
\toprule
File & Contents \\
\midrule
\code{\_\_init\_\_.py} & \code{ssRIRPA(mf, ...)} factory: UHF goes to \code{ssRIURPA}, a non-\code{None}
  \code{rixc} goes to \code{ssRIRRPA}, and everything else goes to \code{ssRIdRRPA}. \\
\code{RIRPA.py} & \code{ssRIRRPA}: the general (xc-kernel-capable) restricted class, which works in the
  \textbf{spin-orbital} $2\ov$ space; also \code{construct\_product\_RI},
  \code{construct\_inverse\_RI} and \code{compress\_low\_rank}. \\
\code{RIdRRPA.py} & \code{ssRIdRRPA}: specialised direct-RPA RHF path in \textbf{spatial} $\ov$
  with explicit spin factors (the path most embedding calls take). \\
\code{RIURPA.py} & \code{ssRIURPA}: unrestricted variant; overrides only the integral/xc construction. \\
\code{NI\_eval.py} & Numerical-integration framework: Clenshaw--Curtis on infinite and
  semi-infinite ranges with a nested error estimate, Gauss--Laguerre, and the optimisation
  of the grid scale $a$. \\
\code{momzero\_NI.py} & Integrands for the zeroth moment (\code{MomzeroDeductNone/D/HigherOrder},
  plus the \code{\_dRHF} variant) and the analytic-offset integrand (\code{MomzeroOffsetCalcGaussLag}). \\
\code{energy\_NI.py} & Integrand for $\Tr (MP)^{1/2}$ (\code{NITrRootMP}, plus \code{\_dRHF}). \\
\bottomrule
\end{tabularx}
\end{center}

\subsection{The mathematics as coded}

This follows Furche (PRB \textbf{64}, 195120) and arXiv:2301.09107. For restricted dRPA in
spatial orbitals,
\begin{align}
  A-B &= D \quad\text{(diagonal)}, &
  A+B &= D + 4L^{\mathsf T}L, \\
  MP &\equiv (A-B)(A+B) = D^2 + S_L^{\mathsf T}S_R \quad(\text{rank }\Naux), &
  \eta_0 &= (MP)^{1/2}(A+B)^{-1}.
\end{align}
Vayesta stores $A+B$ as $(\sqrt2 L)^{\mathsf T}(\sqrt2 L)$ with a further factor of 2 for
spin. The square root is evaluated as a frequency integral,
\begin{equation}
  (MP)^{1/2} = \frac{1}{\pi}\int_{-\infty}^{\infty}
     \Big[\,I - \omega^2 (MP+\omega^2)^{-1}\Big]\,d\omega ,
\end{equation}
and Woodbury reduces the resolvent to auxiliary space. At each frequency,
\begin{align}
  F(\omega) &= (D^2+\omega^2)^{-1} \quad\text{(diagonal)},\\
  Q(\omega) &= S_R\,F(\omega)\,S_L^{\mathsf T} \quad(\Naux\times\Naux)
     \qquad\leftarrow\ \text{the } \Naux^2\ov \text{ step},\\
  \text{term}(\omega) &= (TF)\,S_L^{\mathsf T}\big[(I+Q)^{-1}-I\big]\,S_R F
     \quad(\nt\times\ov).
\end{align}
The default \code{"HO"} scheme subtracts the diagonal and the first-order term $-I$ inside the
integrand and adds them back analytically. The first-order offset,
\begin{equation}
  T\!\int_0^\infty e^{-\omega D}\,S_L^{\mathsf T}S_R\,e^{-\omega D}\,d\omega,
  \qquad\text{exactly}\quad \frac{(S_L^{\mathsf T}S_R)_{pq}}{D_p+D_q},
\end{equation}
would cost $\mathcal{O}(\nt\,\ov^2)$ if done exactly. Vayesta therefore evaluates it with a
\textbf{second, separate 48-point Gauss--Laguerre integral}, whose scale is also optimised.

The remaining steps of the zeroth moment are as follows.
\begin{enumerate}
  \item Fit the grid scale $a$ for each integral (\code{opt\_quadrature\_diag}). The fit uses an
        exponential search over 14 decades, then Newton or bounded minimisation, so that the rule
        integrates a diagonal approximation exactly. Each probe costs $\mathcal{O}(\Np\,\ov)$.
  \item Integrate with Clenshaw--Curtis (\code{npoints=48}, rounded up to a multiple of 4). The
        $\tfrac12$- and $\tfrac14$-grids are nested inside it, and they give the error estimate.
  \item Form the explicit inverse \code{np.linalg.inv(I + Q)} at every point.
  \item Multiply by $(A+B)^{-1}$, itself built as an RI correction (\code{construct\_inverse\_RI}:
        an $\Naux\times\Naux$ inverse plus an SVD; or, in the dRPA path, an explicit
        $\big(I+2LD^{-1}L^{\mathsf T}\big)^{-1}$).
  \item Optionally compress (\code{compress > 0}) by a non-symmetric \code{eig} of
        $S_LS_R^{\mathsf T}$ with an absolute $10^{-12}$ cut. This is off by default.
\end{enumerate}

\paragraph{Higher moments} use the exact recurrence
$\eta_{n+2} = \eta_n D^2 + (\eta_n S_R^{\mathsf T})S_L$. Each order is two GEMMs of
$\mathcal{O}(\nt\Naux\ov)$, which is already optimal.

\paragraph{The energy} (\code{kernel\_energy}) is
$E_c = \tfrac12\big[\Tr(MP)^{1/2} - \Tr A\big]$. The trace is another optimised
Clenshaw--Curtis integral (\code{NITrRootMP}), and its diagonal and first-order parts are
added analytically.

\paragraph{The xc-kernel path} (\code{ssRIRRPA} with \code{rixc}, used by EDMET) differs in
two ways. First, $A-B$ is no longer diagonal, so \code{construct\_product\_RI} builds the
non-symmetric product with rank $\Naux(A{+}B)+\Naux(A{-}B)$, which is up to 2--3$\times$ the
CDERI rank. Second, the whole calculation runs in the $2\ov$ spin-orbital space.

\paragraph{\code{get\_gap}} runs a \textbf{non-symmetric} Davidson
(\code{pyscf.lib.eig}, left and right vectors) on $MP$.

\subsection{Current cost}

Per call of \code{kernel\_moms(0, target\_rot)} on the dRPA RHF path (\code{ssRIdRRPA}),
with $\Np=48$:

\begin{center}
\small
\begin{tabularx}{\textwidth}{@{}>{\raggedright\arraybackslash}p{5.0cm}>{\raggedright\arraybackslash}p{4.4cm}L@{}}
\toprule
Step & Cost & Notes \\
\midrule
CDERI transform \code{get\_cderi} (if \code{lov} not passed) & $\mathcal{O}(\Naux\,n_{\mathrm{ao}}^2\,o)$
  & $N^4$; repeated by every caller that does not pass \code{lov} \\
Build $S_L,S_R$ & $\mathcal{O}(\Naux\ov)$ & two to three $\Naux\times\ov$ arrays held at once \\
Grid fit, both integrals & $\mathcal{O}(N_{\mathrm{it}}\Np\ov)$ & cheap, but runs every call \\
Main integral: $Q(\omega)$ & $\Np\,\Naux^2\ov$ & \textbf{dominant for small $\nt$} \\
Main integral: $\mathrm{inv}(I+Q)$ & $\Np\,\Naux^3$ & explicit inverse \\
Main integral: target contractions & $\Np(2\nt\Naux\ov + \nt\Naux^2)$ & \textbf{dominant for large $\nt$} \\
Offset integral (Gauss--Laguerre) & $\Np\cdot 2\nt\Naux\ov$
  & a second pass, as expensive as the target part of the main integral \\
$(A+B)^{-1}$ & $\Naux^2\ov + \Naux^3 + 2\nt\Naux\ov$
  & needed only because of the $(MP)^{1/2}(A+B)^{-1}$ form \\
Spin stacking (spin-orbital targets) & $\mathcal{O}(\nt^2\ov + \nt^3)$
  & \code{compress\_target\_rot} forms an $\nt\times\nt$ Gram \\
Higher moments & $2\nt\Naux\ov$ per order & fine \\
\bottomrule
\end{tabularx}
\end{center}

\noindent The leading terms are therefore
\begin{equation}
  \text{cost}(\eta_0) \approx 48\,\Naux^2\ov + {\sim}190\,\nt\Naux\ov + 48\,\Naux^3,
  \qquad
  \text{memory} \approx (2\text{--}3)\,\Naux\ov + (2+n_{\mathrm{mom}})\,\nt\ov .
\end{equation}
This is $\mathcal{O}(N^4)$ when $\nt=\mathcal{O}(N)$, and $\mathcal{O}(N^5)$ for the full
moment ($\nt=\ov$, which the code warns about). The energy costs ${\sim}96\,\Naux^2\ov$
($Q$ plus the back contraction at every point). For the xc-kernel path, replace $\Naux$ by
the 2--3$\times$ larger product rank and $\ov$ by $2\ov$.

\paragraph{How the call pattern turns this into total cost.}
\begin{itemize}
  \item \textbf{RPA bath:} a new \code{ssRIdRRPA} is built for every fragment. So the
        fragment-independent $48\,\Naux^2\ov$ work (the $Q(\omega)$ matrices and inverses) and
        the grid fit are repeated $\Nfrag$ times. With a dense Kronecker target of
        $\nt\sim o\,v_{\mathrm{clus}}$ rows, each fragment is $\mathcal{O}(N^4)$, and the
        calculation is $\mathcal{O}(N^5)$ overall.
  \item \textbf{Screening:} one call with all fragments stacked. It forms the full
        $\nt\times2\ov$ result and then keeps only the diagonal fragment blocks
        $T_f\eta_0T_f^{\mathsf T}$, so most of the $\nt\times\ov$ back-contraction is thrown away.
  \item \textbf{Accuracy:} the error is controlled only by the nested-Clenshaw--Curtis
        estimate and a warning above \code{err\_tol}$=10^{-6}$ per element. It is not
        certified. Section~\ref{sec:established} gives \code{momentGW}'s measurement of this same
        route at about $10^{-9}$ relative error on a small-gap system, and an unfitted 48-point
        Clenshaw--Curtis rule measured $5\times10^{-7}$ on a larger one.
\end{itemize}

%==========================================================================
\section{Bringing in the monomial / Davidson improvements}

\subsection{What those projects established}\label{sec:established}

Both \code{monomial\_GW} (a momentGW fork whose legacy Clenshaw--Curtis $\eta_0$ route \emph{is}
this Vayesta route, Gauss--Laguerre offset included) and \code{Davidson-gw} (\code{gwfast})
replaced the frequency integral with the \textbf{Hale--Higham--Trefethen (HHT) / Zolotarev
rational approximation of the inverse square root}. They write it in the symmetric,
$(A+B)^{-1}$-free form
\begin{equation}
  \eta_0 = D^{1/2}\,\Mt^{-1/2}\,D^{1/2},\qquad
  \Mt = D^2 + 4WW^{\mathsf T},\qquad W = D^{1/2}L^{\mathsf T}
  \quad(\text{symmetric positive definite}),
\end{equation}
\begin{equation}
  x^{-1/2} \approx \sum_{j=1}^{\Np}\frac{w_j}{x+s_j}
  \qquad\text{on a certified } [\lambda_{\min},\lambda_{\max}] \supseteq \operatorname{spec}(\Mt).
\end{equation}
Each pole is one resolvent, reduced exactly to auxiliary space by Woodbury:
\begin{align}
  F_j &= D\,(D^2+s_j)^{-1} \quad\text{(diagonal)},\\
  C_j &= I + 4\,L F_j L^{\mathsf T} \quad(\Naux\times\Naux,\ \text{SPD})
       \qquad\leftarrow\ \text{one Gram, } \Naux^2\ov,\\
  \eta_0 &\approx \sum_j w_j\Big[F_j - 4\,F_jL^{\mathsf T}C_j^{-1}LF_j\Big].
\end{align}

\noindent The components, all implemented and tested in at least one of the two repositories:
\begin{enumerate}
  \item \textbf{Certified spectral interval.}
        \[\lambda_{\min}=\min(D)^2,\qquad
          \lambda_{\max}=\max(D^2)+4\min\big(\|W\|_F^2,\ \|W\|_1\|W\|_\infty\big).\]
        \code{gwfast} also has a tighter \code{"spectral"} option, which uses the top
        eigenvalue of one $\Naux\times\Naux$ Gram. Both ends are padded outwards.
  \item \textbf{Pole count from a measured scalar error.} The rule is the smallest $\Np$ for
        which $\max_x|1-\sqrt{x}\,r(x)|\le\mathrm{tol}$ over the interval. $\Np$ grows only like
        $\log\kappa\,\log(1/\mathrm{tol})$. The coefficients come from AGM and
        descending-Landen elliptic functions in extended precision, so $1-k^2$ is never formed
        (\code{momentGW/eta0.py}, \code{gwfast/quadrature.py}).
  \item \textbf{No grid fit and no offset integral.} The rational sum absorbs the bare and
        first-order terms: at zero coupling each solve is the identity. This removes Vayesta's
        \code{"HO"}/Gauss--Laguerre second pass and its $a$-optimisation.
  \item \textbf{Cholesky solves with residual checks instead of \code{inv}.} \code{momentGW}
        measured that this is an accuracy rule rather than a speed rule, because both are
        dominated by the $\Naux^2\ov$ apply.
  \item \textbf{No $\ov\times\ov$ intermediate.} The code records intermediate shapes and a test
        asserts that none is formed.
  \item \textbf{Natural auxiliary functions} (\code{gwfast.integral.natural\_auxiliary\_functions}).
        These rotate the CDERIs to the eigenvectors of $LL^{\mathsf T}$ above a threshold. The
        rotation is symmetric (\code{eigh}) and does not change the result above that threshold.
        Compare Vayesta's off-by-default, non-symmetric \code{eig} compression.
  \item \textbf{Square root and inverse square root together}
        (\code{gwfast.rpa.get\_root\_operators}). The expansion
        \[\Mt^{1/2}\approx\sum_j w_j\big[I - s_j(\Mt+s_j)^{-1}\big]\]
        reuses the same $C_j$, so one pass over the poles gives both.
\end{enumerate}

\noindent Measured results, quoted from those projects:

\begin{center}
\small
\begin{tabularx}{\textwidth}{@{}>{\raggedright\arraybackslash}p{4.3cm}>{\raggedright\arraybackslash}p{4.6cm}L@{}}
\toprule
Source & Comparison & Result \\
\midrule
\code{monomial\_GW/ROADMAP.md}, Milestone 2
  & ozone/PBE (condition $6.7\times10^4$) $\eta_0$ vs a dense oracle
  & Clenshaw--Curtis-48: $1.15\times10^{-9}$; HHT-26: $2.1\times10^{-13}$ \\
same & benzene/cc-pVDZ $\eta_0$ stage & 20 poles vs 48 points, \textbf{1.82$\times$} faster \\
\code{Davidson-gw/docs/ hht-quadrature- benchmark-2026-09-22.md}
  & C$_8$H$_{18}$/def2-TZVPP ($\ov=15411$, $\Naux=1148$)
  & unfitted Clenshaw--Curtis-48: $5\times10^{-7}$; HHT-20: $3\times10^{-14}$;
    \textbf{1.95$\times$} faster at 20 poles, 2.72$\times$ at 16; peak memory roughly halved \\
\bottomrule
\end{tabularx}
\end{center}

\subsection{Applying it to Vayesta's projected moments}

Vayesta needs $T\eta_0$ for a general target $T$ ($\nt\times\ov$), not just
$\eta_0L^{\mathsf T}$ as in GW. The Woodbury form still applies:
\begin{equation}
  T\eta_0 \approx \underbrace{T\,\diag\Big(\sum_j w_jF_j\Big)}_{\approx T}
  \;-\; 4\sum_j w_j\,
  \underbrace{(TF_jL^{\mathsf T})}_{\nt\Naux\ov}\;
  \underbrace{C_j^{-1}}_{\text{solve}}\;
  \underbrace{(LF_j)}_{\nt\Naux\ov}.
\end{equation}
Per pole this costs $\Naux^2\ov$ (Gram) $+\ \Naux^3/3$ (Cholesky)
$+\ 2\nt\Naux\ov + \nt\Naux^2$. The scaling class is the same as today. The gain comes from
having fewer passes and dropping the redundant ones:

\begin{center}
\small
\begin{tabularx}{\textwidth}{@{}>{\raggedright\arraybackslash}p{3.6cm}LLc@{}}
\toprule
& Current (CC-48, \code{"HO"}) & HHT, $\Np\approx20$ & Ratio \\
\midrule
$\Naux^2\ov$ work & 48 & 20 & 2.4$\times$ \\
$\nt\Naux\ov$ work & ${\sim}96$ (main) $+ {\sim}96$ (offset) $+ 2$ & ${\sim}40$ & ${\sim}4.8\times$ \\
$\Naux^3$ & 48 explicit inverses + SVD & 20 Cholesky factorisations & ${\sim}7\times$ \\
Grid fitting, $(A+B)^{-1}$ & yes & not needed & -- \\
Error control & heuristic, about $10^{-7}$ to $10^{-9}$ & certified to tol ($10^{-14}$ default in \code{momentGW}) & -- \\
\bottomrule
\end{tabularx}
\end{center}

The expected $\eta_0$ speed-up is therefore about \textbf{2.4$\times$ when $\nt\ll\Naux$,
about 4$\times$ when $\nt\sim\Naux$, and about 4.8$\times$ when $\nt\gg\Naux$}, with accuracy
improving by 5 to 7 orders of magnitude. These are operation-count estimates and must be
measured. The one directly comparable measurement is 1.82$\times$ on the GW $\eta_0$ stage,
where $\nt=\Naux$.

\noindent The same pole factors give further savings that depend on what each caller actually needs:
\begin{enumerate}
  \item \textbf{Cache the per-pole factors across targets and fragments.} $F_j$, $C_j$ and its
        Cholesky factor depend only on the mean field, not on $T$. Build them once per
        calculation, at $\Np\Naux^2\ov$ and $\Np\Naux^2$ memory, and pass them to every caller.
        For the \textbf{RPA bath} this changes the fragment-independent part from
        $\Nfrag\cdot48\,\Naux^2\ov$ to $20\,\Naux^2\ov$ once. That is a reduction of
        ${\sim}2.4\,\Nfrag\times$ on that term, and it changes the overall scaling from
        $\mathcal{O}(N^5)$ to $\mathcal{O}(N^4)$ plus a per-fragment target term.
  \item \textbf{Two-sided projections.} Several callers contract the $\ov$ index of $\eta_0$
        on both sides, and the left and right targets need not be the same, or even have the
        same number of rows. With a left target $X$ ($n_X\times\ov$) and a right target $Y$
        ($n_Y\times\ov$), $\eta_0$ is symmetric, so
        \begin{equation}
          \begin{aligned}
          X\eta_0Y^{\mathsf T} &\approx X\,\diag\Big(\sum_j w_jF_j\Big)Y^{\mathsf T}
            - 4\sum_j w_j\,B_j^X C_j^{-1}\big(B_j^Y\big)^{\mathsf T},\\
          B_j^X &= XF_jL^{\mathsf T}\ (n_X\times\Naux),\qquad B_j^Y = YF_jL^{\mathsf T}\ (n_Y\times\Naux).
          \end{aligned}
        \end{equation}
        The result carries only target labels ($n_X\times n_Y$), with no particle--hole index
        left. When the two target labels are contracted against each other it is a single
        number, labelled only by the pair $(X,Y)$. That is a trace in auxiliary space:
        \begin{equation}
          \Tr\big[X\eta_0Y^{\mathsf T}\big] \approx
            \Tr\Big[X\,\diag\Big(\sum_j w_jF_j\Big)Y^{\mathsf T}\Big]
            - 4\sum_j w_j\,\Tr\Big[C_j^{-1}\big(B_j^Y\big)^{\mathsf T}B_j^X\Big].
        \end{equation}
        Neither form ever builds the $n_X\times\ov$ intermediate $X\eta_0$. Per pole the cost
        is $(n_X+n_Y)\Naux\ov$ for the two $B$ matrices, which drops to $n_X\Naux\ov$ when
        $Y=X$. Then the open form costs $n_X\Naux^2 + n_Xn_Y\Naux$, and the trace costs
        $\min(n_X,n_Y)\Naux^2+\Naux^3$ instead. Today the cost is $2n_X\Naux\ov$ (main
        integral) plus $2n_X\Naux\ov$ (offset) per point, and then a further $n_Xn_Y\ov$ for
        the projection.

        Where this applies in Vayesta:
        \begin{center}
        \footnotesize
        \begin{tabularx}{\linewidth}{@{}>{\raggedright\arraybackslash}p{3.9cm}>{\raggedright\arraybackslash}p{2.0cm}L>{\raggedright\arraybackslash}p{1.6cm}@{}}
        \toprule
        Caller & $X$ & $Y$ & Form \\
        \midrule
        \code{screening\_moment.py} (\code{dot(momzero[..], rot.T)}) & $T_f$ & $T_f$
          & open, diagonal blocks only \\
        \code{bath/rpa.py} (\code{dot(m0, target\_rot.T)}) & $R_o\otimes R_v$ & the same & open \\
        \code{qemb.py} cumulant energy (\code{m0 - l\_} against \code{l\_2})
          & \code{l\_} & \code{l\_2}, which differs from \code{l\_} by the sign of the negative-CDERI rows & trace \\
        \code{RIRPA.kernel\_energy}, \code{"linear"} xc correction
          & \code{ri\_b\_xc[0]} & \code{ri\_b\_xc[1]} & trace \\
        \code{RIRPA.} \code{direct\_AC\_integration} & \code{lrot} & \code{rrot} & trace \\
        \bottomrule
        \end{tabularx}
        \end{center}

        When either side lies in the RI span the matching $B$ is already known. For $Y=L$,
        $B_j^Y = LF_jL^{\mathsf T} = (C_j-I)/4$. The cumulant energy without negative CDERIs
        has $X=Y=L$, and its trace reduces to
        $-\tfrac14\sum_j w_j\Tr\big[C_j - 2I + C_j^{-1}\big]$ plus the diagonal term. That is
        one Gram and one Cholesky factorisation per pole, with no target contraction at all.

        Keep the rational diagonal $\sum_j w_jF_j$ rather than replacing it by $I$. The callers
        subtract the bare target (\code{m0 - l\_}, \code{eta0 - lrot}), and the consistent
        diagonal makes that difference come out of the low-rank term alone instead of from a
        cancellation of two large numbers.
  \item \textbf{Kronecker-structured targets.} The bath target is $R_o\otimes R_v$, but it is
        stored densely, and \code{screening\_moment} does the same per spin. Applying it in
        factored form,
        $B_j[I,A,Q]=\sum_{ia}R_o[I,i]\,R_v[A,a]\,(F_jL^{\mathsf T})[ia,Q]$ done as two small
        GEMMs, saves a factor of about $o$ (or $v$) per contraction. With item 1, the bath then
        costs about $\mathcal{O}(o\,v\,v_{\mathrm{clus}}\Naux)$ per fragment and pole, which is
        $\mathcal{O}(N^3)$ per fragment and $\mathcal{O}(N^4)$ overall.
  \item \textbf{Targets inside the RI span.} For the cumulant energy in \code{qemb} the target is
        $L$ itself, and $TF_jL^{\mathsf T}$ \emph{is} the Gram, $(C_j-I)/4$. This is exactly the
        \code{momentGW} kernel, at ${\sim}2\Naux^2\ov$ per pole and no $\nt$ term.
  \item \textbf{Energy from the same poles.} Applying the rule to $\sqrt{x}=x\cdot x^{-1/2}$ for
        both $\Mt$ and $D^2$, and differencing so that the leading terms cancel identically, gives
        \begin{equation}
          \Tr\big(\Mt^{1/2}\big) - \Tr(D) \approx
          4\sum_j w_j s_j\,\Tr\!\Big[C_j^{-1}\big(L\,D\,(D^2+s_j)^{-2}L^{\mathsf T}\big)\Big].
        \end{equation}
        This costs ${\sim}2\Naux^2\ov$ per pole, against ${\sim}96\,\Naux^2\ov$ for
        \code{NITrRootMP} today: about 2.4$\times$ faster and certified. One caveat: the bound is
        relative per eigenvalue, so the absolute energy error is at most about
        $2\,\mathrm{tol}\,\Tr(D)$. Choose tol against the required $E_c$ accuracy. At
        $10^{-14}$ this is well under a microhartree for any system Vayesta handles.
  \item \textbf{Natural auxiliary functions.} Rotating $L$ once to $\Naux'\le\Naux$ NAFs saves
        $(\Naux/\Naux')^2$ on every Gram and $\Naux/\Naux'$ on every target contraction. The
        saving depends on the system and threshold and has not been measured here. The rotation
        must use a scale-aware PSD criterion and report the rank and discarded norm, which is
        \code{monomial\_GW} Milestone 4's open item. The current absolute $10^{-12}$ cut does
        neither.
  \item \textbf{Symmetric gap and eigenvectors.} $\Mt$ is symmetric with eigenvalues $\Omega^2$.
        \code{get\_gap} can therefore use the symmetric \code{pyscf.lib.davidson} on
        $x\mapsto D^2x + 4W(W^{\mathsf T}x)$ ($\mathcal{O}(\Naux\ov)$ per matvec) instead of the
        left/right non-symmetric \code{eig} on $MP$ followed by re-biorthogonalisation.
\end{enumerate}

\noindent Rough projected totals for the two main embedding uses (operation counts, before NAF):

\begin{center}
\small
\begin{tabularx}{\textwidth}{@{}>{\raggedright\arraybackslash}p{2.4cm}LL>{\raggedright\arraybackslash}p{3.0cm}@{}}
\toprule
Use & Today & With the changes & Gain \\
\midrule
RPA bath, $\Nfrag$ fragments
  & $\Nfrag\big(48\Naux^2\ov + {\sim}190\,o\,v_c\Naux\ov\big)\sim\mathcal{O}(N^5)$
  & $20\Naux^2\ov + \Nfrag\cdot20\big(o\,v\,v_c\Naux + \dots\big)\sim\mathcal{O}(N^4)$
  & grows with system size; already about $\Nfrag\times o$ on the target term \\
\code{mrpa} screening
  & $48\Naux^2\ov + {\sim}190\,\nt\Naux\ov$, \ $\nt=\sum_f2\ov_f$
  & $20\Naux^2\ov + 20\sum_f\big(2n_f\Naux\ov/k + n_f\Naux^2\big)$, where $k$ is the Kronecker saving
  & about 5$\times$, and more with Kronecker targets \\
Total dRPA energy & ${\sim}96\Naux^2\ov$ & ${\sim}40\Naux^2\ov$ & about 2.4$\times$, certified \\
\bottomrule
\end{tabularx}
\end{center}

\subsection{Proposed sequence}

This mirrors \code{monomial\_GW}'s ordering: restricted dRPA first, measured before it becomes
the default, and the old route kept as an independent reference.
\begin{enumerate}
  \item \textbf{Baseline.} Record $\eta_0$, the projected moments, energies, screened
        interactions and bath occupations for a small set (H$_2$O, a small-gap molecule, one
        Hubbard case, one solid), using the existing \code{vayesta/tests/rpa/test\_rirpa\_*}
        systems plus a dense \code{ssRPA} oracle.
  \item \textbf{HHT kernel for \code{ssRIdRRPA}.} Port \code{momentGW/eta0.py} (interval, poles,
        scalar certificate) and a projected apply that accepts $T$ in one-sided, two-sided and
        in-RI-span forms. Add an \code{eta0\_method=\{"clencur","hht"\}} option. Keep
        Clenshaw--Curtis as the default until the baseline agrees within the certified error,
        then switch.
  \item \textbf{Pole-factor cache and caller changes.} Build the factors once per embedding
        object. Move \code{bath/rpa.py} and \code{screening\_moment.py} to the two-sided,
        Kronecker-factored apply, and stop recomputing CDERIs per fragment by passing
        \code{lov} everywhere.
  \item \textbf{Energy and gap} from the same poles (items 5 and 7 above).
  \item \textbf{NAF/scale-aware compression} behind an option, reporting rank, discarded norm
        and the observed effect on $\eta_0$ and energy.
  \item \textbf{Deferred: the UHF and xc-kernel paths.} \code{ssRIURPA} should follow with the
        stacked-spin convention that \code{gwfast} already validated (spin factor 1 over
        $2\ov$). EDMET's \code{rixc} path cannot use the symmetric
        $D^{1/2}\Mt^{-1/2}D^{1/2}$ form, because $A-B$ is not diagonal. It needs either a
        non-symmetric Woodbury rational apply (LU solves) on $MP$, whose interval is not
        theorem-backed, or $(A-B)^{1/2}$ as a diagonal-plus-low-rank function. Keep it on
        Clenshaw--Curtis until that is designed. Periodic \code{kdf} cases also wait, for the
        reasons \code{monomial\_GW} Milestone 6 records.
\end{enumerate}

\subsection{Points to settle before implementation}

\begin{itemize}
  \item \textbf{Spin conventions.} \code{ssRIdRRPA} mixes spatial integrals with spin-orbital
        outputs (\code{stack\_spin}, \code{return\_spatial}). The factor 4 in $\Mt$ and the split
        of the spin-diagonal identity term must be checked against the dense \code{ssRPA} before
        anything else.
  \item \textbf{Negative CDERIs} (\code{lov\_neg}, from non-PSD metrics in some periodic setups).
        $A+B$ then takes an indefinite low-rank part, $C_j$ is no longer SPD, and the certified
        interval argument fails. The HHT path should refuse them at first.
  \item \textbf{Memory for the cache.} Storing the factors costs $\Np\Naux^2$ (about 0.2\,GB at
        $\Naux=1000$, $\Np=20$). For very large $\Naux$, rebuild per pole; \code{gwfast}'s
        \code{PRECOMPUTE\_S} switch is exactly this trade.
  \item \textbf{Higher moments} are already exact recurrences at $\mathcal{O}(\nt\Naux\ov)$ per
        order. \code{monomial\_GW} Milestone 3 found that its high-order instabilities were in
        the downstream realisation, not in this recurrence, so no change is proposed there.
\end{itemize}

%==========================================================================
\section{Local and complementary projections (draft)}\label{sec:local}

\noindent\textbf{Status: first draft for discussion.} All costs are operation counts. Nothing
here has been implemented or measured. It builds on the pole form of
Section~\ref{sec:established} and on the discussion of local spaces labelled by $x$.

\subsection{Setup}

Each label $x$ defines a local occupied and a local virtual space through orthonormal rows of
$C^x$:
\begin{equation}
\begin{aligned}
  &c_o\ (o_x\times o),\quad c_v\ (v_x\times v), && c_oc_o^{\mathsf T}=I,\quad c_vc_v^{\mathsf T}=I,\\
  &P_o=c_o^{\mathsf T}c_o,\quad P_v=c_v^{\mathsf T}c_v, && Q_o=I-P_o,\quad Q_v=I-P_v,\\
  &e_o,\ e_v\ \text{complement rows}, && e_o^{\mathsf T}e_o=Q_o,\quad e_v^{\mathsf T}e_v=Q_v,\\
  &X=c_o\otimes c_v && \text{local ph target, } n_x=o_xv_x \text{ rows}.
\end{aligned}
\end{equation}
$n_x$ is $\mathcal{O}(1)$ in system size. There are $\mathcal{O}(N)$ labels for molecules and
$\mathcal{O}(1)$ for solids.

The right-hand side of the quantities discussed is a \textbf{complement} of the local space.
Two complements have come up, and they are different quantities, not two representations of
one:

\begin{center}
\small
\begin{tabularx}{\textwidth}{@{}>{\raggedright\arraybackslash}p{2.3cm}>{\raggedright\arraybackslash}p{3.0cm}L>{\raggedright\arraybackslash}p{3.3cm}@{}}
\toprule
Name & Projector on the right & Blocks kept & Example \\
\midrule
ph complement & $\bar\Pi = I - P_o\otimes P_v$
  & everything except local--local, including the mixed local-occ/env-vir and env-occ/local-vir blocks
  & $\tilde\eta^x_{\bar i\bar a,kc}$ (full MO indices on the right) \\
strict environment & $Q_o\otimes Q_v$ & env--env only
  & $\eta^x_{\bar i\bar a,\bar k\bar c}$ (complement barred indices) \\
\bottomrule
\end{tabularx}
\end{center}

Either can be \textbf{represented} in the full MO basis (apply the projector) or in the
complement basis (rotate with $e_o\otimes e_v$). As George notes, the basis is a choice.
Section~\ref{sec:kernelA} shows that the MO-basis representation is the cheaper one to
compute.

\noindent\emph{To decide:} which complement each use needs.

\subsection{Shared and per-label factors}

Shared factors, built once per mean field, as in Section~\ref{sec:established}:
\begin{equation}
  F_j = D(D^2+s_j)^{-1},\qquad C_j = I+4LF_jL^{\mathsf T}\ \to\ \mathrm{chol}(C_j),\qquad
  \eta_0 \approx \sum_j w_j\big[F_j - 4F_jL^{\mathsf T}C_j^{-1}LF_j\big].
\end{equation}
They cost $\Np(\Naux^2\ov+\Naux^3/3)$ and need $\Np\Naux^2$ memory.

Per-label left factors, built for each $x$ and each pole:
\begin{align}
  B_j^x &= XF_jL^{\mathsf T}\ \ (n_x\times\Naux), &&\text{Kronecker: } (o_x\,o\,v + o_xv_x\,v)\,\Naux,\\
  G_j^x &= w_j\,B_j^xC_j^{-1}\ \ (n_x\times\Naux), &&\text{one Cholesky solve: } n_x\Naux^2 .
\end{align}
$B_j^x$ is applied in factored form: reshape $F_jL^{\mathsf T}$ to $(o,v,\Naux)$ and contract
$c_o$, then $c_v$. It is never formed as a dense $n_x\times\ov$ target.

Every quantity below is built from these factors, through one of three kernels.

\subsection{Kernel A: one-sided, with a projector on the right}\label{sec:kernelA}

\begin{equation}
  X\eta_0R,\qquad R\in\{\,I,\ \bar\Pi,\ Q_o\otimes Q_v\,\}.
\end{equation}

\paragraph{Bare term.} It is $XD_rR$ with $D_r=\diag(\sum_jw_jF_j)\approx I$. For $R=I$ it is
about $X$. For either complement $XR=0$ exactly, so the bare term is $X(D_r-I)R$, which is
$\mathcal{O}(\mathrm{tol})$. It can be dropped, leaving only the correlated part, with no large
cancellation.

\paragraph{Low-rank term.} Sum over the poles \emph{before} projecting:
\begin{equation}
  S^x = -4\sum_j G_j^x\,(LF_j)\quad(n_x\times\ov,\ \ \Np\,n_x\Naux\ov),\qquad
  X\eta_0R \approx S^xR .
\end{equation}
The projector has to come after the pole sum. $F_j$ is diagonal in the MO basis and does not
commute with $R$, so applying $R$ inside the sum would cost one application per pole.
Here $R$ is applied once per label:

\begin{center}
\small
\begin{tabularx}{\textwidth}{@{}>{\raggedright\arraybackslash}p{2.9cm}L>{\raggedright\arraybackslash}p{4.4cm}@{}}
\toprule
$R$ & How & Cost per label \\
\midrule
$\bar\Pi$ & $S-(SX^{\mathsf T})X$, where $SX^{\mathsf T}=-4\sum_jG_j^x(B_j^x)^{\mathsf T}$ is the two-sided
  block (kernel B, case a), so it is already available
  & $\Np n_x^2\Naux$ + ${\sim}n_x\ov\,o_x$ to expand \\
$Q_o\otimes Q_v$, MO basis & $S - S(P_o\otimes I) - S(I\otimes P_v) + S(P_o\otimes P_v)$, each a low-rank subtraction
  & ${\sim}2\,n_x\ov(o_x+v_x)$ \\
$Q_o\otimes Q_v$, complement basis & rotate $S$ by $e_o^{\mathsf T}$, $e_v^{\mathsf T}$
  & $n_x(o_{\mathrm{env}}ov + o_{\mathrm{env}}v_{\mathrm{env}}v)$; larger, so use only if that basis is needed downstream \\
\bottomrule
\end{tabularx}
\end{center}

\paragraph{Cost per label.}
\begin{equation}
  \Np\Big[(o_x\,o\,v+o_xv_x\,v)\Naux + n_x\Naux^2 + \underbrace{n_x\Naux\ov}_{\text{dominant, output-bound}}\Big]
  + \mathcal{O}\big(n_x\ov(o_x+v_x)\big).
\end{equation}
The output has $n_x\ov$ entries, so the $n_x\Naux\ov$ back-multiply per pole cannot be
avoided while the result is kept open on the right. The only levers on it are $\Np$ and
$\Naux$ (natural auxiliary functions).

\paragraph{Batching across labels.} Stack $G_j^x$ over a batch of labels, so that each pole is
one GEMM $(\sum_xn_x\times\Naux)\cdot(\Naux\times\ov)$. The flops are the same but the efficiency
is better. Consume or store each $S^xR$ before the next batch, because holding all labels is
$\mathcal{O}(N^3)$ memory for molecules.

\subsection{Kernel B: two-sided, $X\eta_0Y^{\mathsf T}$ with a small $Y$}

\begin{equation}
  X\eta_0Y^{\mathsf T} \approx XD_rY^{\mathsf T} - 4\sum_j G_j^x\,(B_j^y)^{\mathsf T},
  \qquad B_j^y = YF_jL^{\mathsf T}\ \ (n_Y\times\Naux).
\end{equation}
Per pole this costs forming $B_j^y$ plus $n_x\Naux n_Y$. The open $n_x\times\ov$ block is
never built. Kernel B pays off only when $n_Y\ll\ov$, or when $B_j^y$ is known without a sweep
over $\ov$. That is the case when the right-hand index is contracted with something small
\emph{before} anything is stored.

\begin{center}
\small
\begin{tabularx}{\textwidth}{@{}>{\raggedright\arraybackslash}p{1.9cm}>{\raggedright\arraybackslash}p{1.8cm}>{\raggedright\arraybackslash}p{3.9cm}>{\raggedright\arraybackslash}p{3.2cm}L@{}}
\toprule
Case & $Y$ & $B_j^y$ & Extra cost per pole & Uses \\
\midrule
(a) same space & $X$ & $B_j^x$ (already built) & $n_x^2\Naux$
  & screening $T_f\eta_0T_f^{\mathsf T}$, RPA bath $m_0T^{\mathsf T}$, the $\bar\Pi$ correction in kernel A \\
(b) RI span & $L$ & $LF_jL^{\mathsf T}=(C_j-I)/4$, free
  & none: $X\eta_0L^{\mathsf T}\approx XD_rL^{\mathsf T}-\sum_jw_jB_j^x(I-C_j^{-1})$
  & GW-type $\eta_0V$; interactions built from RI vectors \\
(c) RI span, projected & $L\bar\Pi$ & $(C_j-I)/4-(LX^{\mathsf T})B_j^x$
  & $\Naux^2n_x$; $LX^{\mathsf T}$ once per label
  & RHS of $\tilde\eta^x$ contracted directly with $(kc|Q)$ \\
(c$'$) RI span, env--env & $L(Q_o\otimes Q_v)$
  & $(C_j-I)/4$ minus the $P_o$, $P_v$ and $P_o\otimes P_v$ corrections, each a factored low-rank product
  & ${\sim}(o_x+v_x)\ov\Naux + (o_xv+ov_x)\Naux^2$
  & RHS of $\eta^x_{\bar i\bar a,\bar k\bar c}$ contracted with integrals \\
(d) another label & $X'=c_o'\otimes c_v'$ & $B_j^{x'}$, already built for $x'$ & $n_x\Naux n_{x'}$
  & cross-label couplings, if needed \\
\bottomrule
\end{tabularx}
\end{center}

For (c) and (c$'$) the bare term $XD_rRL^{\mathsf T}$ is again $\mathcal{O}(\mathrm{tol})$ and
can be dropped.

Whether kernel B beats kernel A for (c$'$) depends on the sizes:
$\Naux(o_xv+ov_x)$ against $o_xv_x\,ov$ per pole. For small local spaces in a large system,
kernel B is expected to win, but this has to be measured.

\subsection{Kernel C: fully contracted, $\Tr[X\eta_0Y^{\mathsf T}]$}

\begin{equation}
  \Tr\big[X\eta_0Y^{\mathsf T}\big] \approx \Tr\big[XD_rY^{\mathsf T}\big]
   - 4\sum_jw_j\Tr\big[C_j^{-1}(B_j^y)^{\mathsf T}B_j^x\big].
\end{equation}
This costs $\min(n_X,n_Y)\Naux^2+\Naux^3$ per pole once the $B$ matrices are built. When
$X=Y=L$, it is $-\tfrac14\sum_jw_j\Tr[C_j-2I+C_j^{-1}]$ plus the diagonal term: one Gram and
one Cholesky factorisation per pole, with no target contraction. The uses are the cumulant
energy, the \code{"linear"} xc energy correction and \code{direct\_AC\_integration}
(Section~3.2, item~2).

\subsection{Summary by use}

\begin{center}
\small
\begin{tabularx}{\textwidth}{@{}>{\raggedright\arraybackslash}p{3.2cm}>{\raggedright\arraybackslash}p{2.6cm}>{\raggedright\arraybackslash}p{1.6cm}>{\raggedright\arraybackslash}p{1.9cm}L@{}}
\toprule
Use & Right side & Output & Kernel & Dominant cost per label \\
\midrule
$\tilde\eta^x_{\bar i\bar a,kc}$ & $\bar\Pi$, MO basis & $n_x\times\ov$ & A & $\Np n_x\Naux\ov$ \\
$\eta^x_{\bar i\bar a,\bar k\bar c}$, stored & $Q_o\otimes Q_v$ & $n_x\times n_{\mathrm{env}}$
  & A, then rotate if the env basis is needed & $\Np n_x\Naux\ov$ \\
either, RHS contracted with $L$ next & $L\bar\Pi$ or $L(Q_o\otimes Q_v)$ & $n_x\times\Naux$ & B (c)/(c$'$)
  & $\Np\Naux^2(o_xv+ov_x)$; no $n_x\ov$ term \\
screening, RPA bath & $T_f$ & $n_x\times n_x$ & B (a) & $\Np n_x\Naux^2$ + factored forward \\
energies & various & scalar & C & $\Naux^2$ per pole \\
\bottomrule
\end{tabularx}
\end{center}

With the shared factors, the totals for kernel A are $\mathcal{O}(N^4)$ for molecules. For
solids they are a shared $\mathcal{O}(N^4)$ term plus $\mathcal{O}(N^3)$ per label. Today's route
costs about $48\Naux^2\ov + 192\sum_xn_x\Naux\ov$ for any of these, and always builds the
$n_x\times\ov$ block. With the shared factors, kernel A costs about
$20\Naux^2\ov + 20\sum_xn_x\Naux\ov$, about 2.4$\times$ on the shared term and about
10$\times$ on the per-label term. Kernels B and C remove the per-label $\ov$ sweep where the
use allows it.

\subsection{Proposed kernel interface}

\begin{small}
\begin{verbatim}
# s_j, w_j, F_j, chol(C_j); certified interval
poles = build_pole_factors(D, L, tol)
# B_j^x, G_j^x (Kronecker-factored)
left = local_left_factors(poles, c_o, c_v)
# kernel A
one_sided(poles, left, right=None | "ph_complement" | "env", basis="mo" | "env")
# kernel B; right = left' | "ri" | ("ri", projector)
two_sided(poles, left, right)
# kernel C
trace(poles, left, right)
\end{verbatim}
\end{small}
Every function also takes a batch of labels, and each records the per-pole Cholesky residuals
and the intermediate shapes, following \code{monomial\_GW} Milestone~2.

\subsection{Open questions}

\begin{enumerate}
  \item \textbf{Which complement} ($\bar\Pi$ or $Q_o\otimes Q_v$) each use needs.
  \item \textbf{What the right-hand index is contracted with next} in each use. This decides
        between kernel A and kernel B, and it is the main source of savings beyond the constant
        factors.
  \item \textbf{Whether the outputs are stored or consumed} per label. This sets the memory plan
        for kernel A.
  \item \textbf{Orthonormality of $C^x$ rows.} The $XR=0$ cancellation and the cheap projector
        applications rely on it.
  \item \textbf{Spin.} Restricted spatial forms are given here, with the factor 4 still to be
        checked against the dense \code{ssRPA}. The spin-orbital stacking used by the existing
        callers has to be mapped onto them.
  \item \textbf{Dropping the $\mathcal{O}(\mathrm{tol})$ bare term.} It is consistent with the
        certificate on $\eta_0$, but the error bound for the projected quantity should be stated
        explicitly, for example by Cauchy--Schwarz with $\|X\|=1$.
\end{enumerate}

\end{document}
